% Part: first-order-logic
% Chapter: proof-systems
% Section: axiomatic-deduction

\documentclass[../../../include/open-logic-section]{subfiles}

\begin{document}

\iftag{FOL}
      {\olfileid{fol}{prf}{axd}}
      {\olfileid{pl}{prf}{axd}}

\olsection{公理推导}

公理推导是最古老、最简单的逻辑推导系统。该系统中的推导仅仅是简单的语句序列。一个语句序列构成一个正确的推导，当且仅当其中的每个语句~$!A$ 都满足以下条件之一：

\begin{enumerate}
\item $!A$ 是一条公理，或
\item $!A$ 是某个给定的语句集合~$\Gamma$ 中的一个元素，或
\item $!A$ 由一条推理规则证成。
\end{enumerate}

作为公理，$!A$ 必须具有某种固定语句模式的形式。存在许多公理模式集合，它们都能为一阶逻辑提供令人满意（即可靠且完备）的推导系统。有些是按照它们所涉及的联结词来组织的，例如，模式
\[
!A \lif (!B \lif !A) \qquad !B \lif (!B \lor !C) \qquad (!B \land !C) \lif !B
\]
就是关于 $\lif$、$\lor$ 和~$\land$ 的常见公理。有些公理系统追求最少数量的公理。根据所选取的原始联结词，甚至有可能找到只包含一条公理的公理系统。

一条推理规则是一个条件句，它为推导中语句的证成提供充分条件。肯定前件是一条非常常见的此类规则：它说的是，如果 $!A$ 和 $!A \lif !B$ 已经被证成，那么 $!B$ 就被证成。这意味着，只要某个语句~$!A$ 和 $!A \lif !B$ 在推导中位于 $!B$ 之前，推导中包含语句~$!B$ 的这一行就是得到证成的。

基于公理推导的 $\Proves$ 关系定义如下：$\Gamma \Proves !A$ 当且仅当存在一个以语句~$!A$ 作为其最后一个公式的推导（其中 $\Gamma$ 取作该推导中由上述(2)证成的语句组成的集合）。若 $!A$ 有一个 $\Gamma$ 为空的推导，即推导中的每一个语句要么由(1)证成，要么由(3)证成，则 $!A$ 是一个定理。例如，下面是一个证明 $\Proves
!A  \lif (!B \lif (!B \lor !A))$ 的推导：

\begin{derivation}
  1. & $!B \lif (!B \lor !A)$ \\
  2. & $(!B \lif (!B \lor !A)) \lif (!A  \lif (!B \lif (!B \lor !A)))$\\
  3. & $!A  \lif (!B \lif (!B \lor !A))$
\end{derivation}

第1行的语句具有公理 $!A \lif (!A \lor !B)$ 的形式（这里 $!A$ 和 $!B$ 的角色互换了）。第2行的语句具有公理 $!A \lif (!B \lif !A)$ 的形式。因此，这两行都是得到证成的。第3行由肯定前件证成：如果我们将其缩写为 $!D$，那么第2行具有 $!C \lif !D$ 的形式，其中 $!C$ 是 $!B \lif (!B \lor !A)$，也就是第1行。

如果 $\Gamma \Proves \lfalse$，则集合 $\Gamma$ 是不一致的。一个完备的公理系统也能证明 $\lfalse \lif !A$ 对任意~$!A$ 成立，因此，如果 $\Gamma$ 不一致，那么 $\Gamma \Proves !A$ 对任意~$!A$ 都成立。

逻辑的公理化推导系统最早由 Gottlob Frege（弗雷格）在其1879年的《Begriffsschrift》中给出，正因如此，该书常被认为是现代逻辑的开山之作。它们后来在 Alfred North Whitehead（怀特海）和 Bertrand Russell（罗素）的《Principia Mathematica》以及 David Hilbert（希尔伯特）与其学生在20世纪20年代的工作中得以完善。因此，它们常被称为``Frege 系统''或``Hilbert 系统''。它们非常通用，因为通常很容易为一种逻辑找到一个公理系统。由于推导的结构非常简单，且只有一两条推理规则，相对而言也容易证明一些\emph{关于}它们的性质。然而，它们很难在实际中使用，也就是说，很难找到并写出证明。

\end{document}